\section{Estimator Details and the Learned Lifting}
\label{app:A1_estimator}

The reported instance is fit by the least-squares problem of
Section~\ref{sec:learning}, on the transition tuples
$\mathcal{D}=\{(s_t,c_t,s_{t+1})\}$ constructed there. This appendix covers
the learned-lifting variant instead, in which the encoder and the decoder are
not the identity. What a lifting is, and why the reported instance leaves it
off, is said in Section~\ref{sec:model}.

\textbf{Training transition.} Fitting the operator needs a one-step
latent prediction from an encoded state and the applied command. Let
\begin{equation}
    q_t
    =
    K E_{\phi}(s_t)+B c_t+b.
    \label{eq:training_transition}
\end{equation}
Here $q_t$ is the predicted encoding of $s_{t+1}$ under the current
parameters, computed from the encoder $E_{\phi}$, the transition $K$,
the command column $B$, and the bias $b$.

\textbf{Training objective.} Three losses jointly shape the encoder,
decoder, and transition. We optimize
\begin{equation}
    \begin{aligned}
        \mathcal{L}_{\mathrm{rec}}
        &=
        \mathbb{E}_{\mathcal{D}}\!\left[
            \left\|
                s_t-D_{\psi}(E_{\phi}(s_t))
            \right\|_2^2
        \right],\\
        \mathcal{L}_{\mathrm{lin}}
        &=
        \mathbb{E}_{\mathcal{D}}\!\left[
            \left\|
                E_{\phi}(s_{t+1})-q_t
            \right\|_2^2
        \right],\\
        \mathcal{L}_{\mathrm{pred}}
        &=
        \mathbb{E}_{\mathcal{D}}\!\left[
            \left\|
                s_{t+1}-D_{\psi}(q_t)
            \right\|_2^2
        \right],\\
        \mathcal{L}
        &=
        \lambda_{\mathrm{rec}}\mathcal{L}_{\mathrm{rec}}
        +
        \lambda_{\mathrm{lin}}\mathcal{L}_{\mathrm{lin}}
        +
        \lambda_{\mathrm{pred}}\mathcal{L}_{\mathrm{pred}}.
    \end{aligned}
    \label{eq:training_objective}
\end{equation}
The nonnegative weights $\lambda_{\mathrm{rec}}$,
$\lambda_{\mathrm{lin}}$ and $\lambda_{\mathrm{pred}}$ set how much
each term contributes to the total.
The reconstruction term $\mathcal{L}_{\mathrm{rec}}$ preserves
information present in the memory state $s_t$. The latent-consistency
term $\mathcal{L}_{\mathrm{lin}}$ pulls the encoded successor toward the
affine prediction $q_t$. The prediction term $\mathcal{L}_{\mathrm{pred}}$
ties that latent transition back to accuracy in the observed state
space.

These three terms are complementary. Latent consistency alone admits
uninformative constant encodings, since a constant encoder can drive
$\mathcal{L}_{\mathrm{lin}}$ to zero without capturing any structure.
Reconstruction alone does not require the representation to support
prediction, since $D_{\psi}(E_{\phi}(s_t))$ can recover $s_t$ while
carrying no information about $s_{t+1}$. Training updates
$(\phi,\psi,K,B,b)$, and the frozen language model's parameters
$\theta$ never change.
